\originalchapter{8.333复习包1：界面热力学与概率}
\osource{8.333 Fall 2013课程页公开的Fall 2007 Test1 Review Problems及官方Solutions；9题，非2013年实际测验卷。采用官方热量定义，补充近似条件。}
\begin{exercise}[原题1(a)--(c)]\label{r:333-1-1}
表面张力$\sigma$以功$\delta W=\sigma\dd A$定义。(a)球形水滴半径$R$的内外压差及肥皂泡压差；(b)空气、固体、水三界面张力$\sigma_{aw},\sigma_{sw},\sigma_{sa}$给接触角与完全润湿条件；(c)室温水$\sigma_0\simeq0.07$ N/m，估重力/毛细分界尺度并举例。
\end{exercise}
\begin{solution}
水滴$\Delta p\,4\pi R^2\dd R=\sigma\,8\pi R\dd R$，故$\Delta p=2\sigma/R$。肥皂泡有内外两张界面，薄膜近似$4\sigma/R$。三相线平行于固面的张力平衡$\sigma_{sa}=\sigma_{sw}+\sigma_{aw}\cos\theta$（$\theta$从水内测），$\cos\theta=(\sigma_{sa}-\sigma_{sw})/\sigma_{aw}$。若$\sigma_{sa}\ge\sigma_{sw}+\sigma_{aw}$，以水膜替代干界面不增加自由能，完全润湿；粗糙、滞后、钉扎可改变观察角，不包含于此平衡模型。

比较$\rho g\ell$与$\sigma/\ell$，毛细长$\ell_c=\sqrt{\sigma/(\rho g)}\simeq2.7$ mm。毫米小液滴趋球形，大水面受重力平整；毛细管弯月面、细针或水黾在水面承重用到这一竞争。
\end{solution}
\begin{exercise}[原题2(a)--(d)]\label{r:333-1-2}
表面活性剂$N$分子铺于面积$A$。(a)解释$\sigma\simeq\sigma_0-N\kb T/A$；(b)肥皂触滴表面时接触角变化；(c)给$\sigma_A=N\kb T/(A-Nb)^2-2aN^2/A^3$、$\sigma_T=-N\kb/(A-Nb)$，求$\sigma(A,T)$及$a,b$意义；(d)用$E_A,\sigma,\sigma_T,\sigma_A$表达定张力与定面积热容差。此处热容为$\delta Q/\dd T$，原题紧接定义的抄写易漏机械功，应遵循官方解的热量式。
\end{exercise}
\begin{solution}
二维理想分子铺展压为$N\kb T/A$，自由能随面积增大降低，张力减少该量。若固体两界面张力近似固定，降$\sigma_{aw}$使前题$\cos\theta$增大、滴更铺开；肥皂同时改变固--水界面时应重新比较全部三张力。

积分两个偏导并用稀薄极限定常数得$\sigma=\sigma_0-N\kb T/(A-Nb)+aN^2/A^2$。$b$为排除面积，吸引$a>0$减少分子铺展压、回升张力。原两偏导交叉相等提供可积性检查。

$\delta Q=\dd E-\sigma\dd A=(E_A-\sigma)\dd A+C_A\dd T$，定$\sigma$有$A_T=-\sigma_T/\sigma_A$，所以
\begin{equation}
C_\sigma-C_A=-(E_A-\sigma)\frac{\sigma_T}{\sigma_A}=T\frac{\sigma_T^2}{\sigma_A},
\end{equation}
最后用$E_A=\sigma-T\sigma_T$，来自$\dd F=-S\dd T+\sigma\dd A$。若只求$\partial_TE|_\sigma$，则还差机械功$\sigma A_T$，不能将这导数误当定张力吸热热容。
\end{solution}
\begin{exercise}[原题3(a)--(c)]\label{r:333-1-3}
理想气温标$\Theta$满足$PV=N\kb\Theta$且$E=E(\Theta)$，不能假设$E\propto\Theta$。做Carnot循环：(a)两等温吸放热；(b)任意热容时的绝热体积变化；(c)证明气体温标与热力学温标等价。
\end{exercise}
\begin{solution}
热端等温$Q_H=N\kb\Theta_H\log(V_2/V_1)$，冷端放热绝对值$Q_C=N\kb\Theta_C\log(V_3/V_4)$，因为等温$E$不变。绝热满足$E'(\Theta)\dd\Theta=-N\kb\Theta\dd V/V$，故$\log(V_f/V_i)=-\int_{\Theta_i}^{\Theta_f}E'(\Theta)\dd\Theta/(N\kb\Theta)$。两条反向绝热线的积分相反，得到$V_2/V_1=V_3/V_4$，因此$Q_H/Q_C=\Theta_H/\Theta_C$。Carnot热力学定义$Q_H/Q_C=T_H/T_C$，温标比例固定，以同一参考点校准即$T=\Theta$，全程没有用常数热容。
\end{solution}
\begin{exercise}[原题4(a)--(c)]\label{r:333-1-4}
(a)由$\dd E=T\dd S-P\dd V$证明理想气内能仅依赖温度；(b)给与$E=E(T)$相容的最一般状态方程；(c)证明van der Waals气$C_V$只依赖温度。
\end{exercise}
\begin{solution}
$F$的Maxwell关系$S_V=P_T$，故$E_V|_T=TS_V-P=TP_T-P$。理想$P=N\kb T/V$给0。一般$E_V=0$等价$\partial_T(P/T)=0$，所以$P=T f(V,N)$，并非只允许理想气；稳定、广延性可继续限制$f$。van der Waals$P=N\kb T/(V-Nb)-aN^2/V^2$给$E_V=aN^2/V^2$，积分$E=E_0(T)-aN^2/V$，因此$C_V=E_0'(T)$，只在采用这套温度无关$a,b$的单相状态方程模型中成立。
\end{solution}
\begin{exercise}[原题5(a)--(f)]\label{r:333-1-5}
一摩尔水在共存$(P,T)$吸收潜热$L$变汽，体积差$\Delta V$，绝热降至$P-\dd P$、$T-\dd T$后凝结。(a)证明$W=\Delta V\dd P+O(\dd P^2)$及CC式；(b)指出以$L(T)-L(T-\dd T)=W$推$L\propto T$的错误；(c)恒潜热、理想汽时$P(T)$；(d)飓风热端80°F、冷端$-120$°F，暖水层至少200英尺，每小时凝结9000万吨，潜热$2.3\times10^6$ J/kg，估最高效率与功率；(e)等温理想空气静力$P(h)$；(f)估9 km珠峰水沸点，列出所用空气温度和分子质量。
\end{exercise}
\begin{solution}
微小Carnot在$P,V$图的面积为$\Delta V\dd P$，绝热边的横移只影响二阶；$W/Q_H=\dd T/T$给$\dd P/\dd T=L/(T\Delta V)$。绝热膨胀本身可凝结部分汽，低温等温段不再把完整一摩尔汽凝结；直接以两端同一摩尔潜热相减遗漏这一阶分率，官方解正指出此点。

取汽$\Delta V\simeq RT/P$，$\log[P(T)/P(T_0)]=-(L/R)(1/T-1/T_0)$。飓风给温$299.8,188.7$ K，最高效率约0.371；质量率$9\times10^{10}/3600=2.5\times10^7$ kg/s，供热$5.75\times10^{13}$ W，最大功约$2.13\times10^{13}$ W。200英尺约61 m暖层是持续蒸发的储库条件，并非此瞬时Carnot功率式的乘数；实际远低于可逆上界。

空气静力$\dd P=-\rho g\dd h$、$\rho=mP/(\kb T_a)$，所以$P(h)=P_0\e^{-mgh/(\kb T_a)}$。取等温$T_a=273$ K、$m=29$原子质量单位，尺度约8.0 km，9 km处约0.324大气压。每质量潜热用$R_v=\kb/m_{H_2O}\simeq461.5$ J kg$^{-1}$ K$^{-1}$，$1/T_b=1/373.15-(R_v/L_v)\log(P/P_0)$，约344 K、71°C。这是题设等温大气估计，不使用现实天气；换空气温度会改变估值。
\end{solution}
\begin{exercise}[原题6(a)--(f)]\label{r:333-1-6}
石英慢冷结晶、放潜热$L$，快冷成玻璃，均近似不可压缩，$\dd E=T\dd S+\mu\dd N$。(a)广延性给$\mu$；(b)晶体$C_X=\alpha T^3$、玻璃$C_G=\beta T$，假设都$S(0)=0$求熵；(c)再假设同$E_0$求内能；(d)化学平衡形式给$T_m$；(e)潜热；(f)检验符号及最可疑假设。这里$\alpha,\beta$按整个样品热容计。
\end{exercise}
\begin{solution}
Euler$E=TS+\mu N$，所以$\mu=(E-TS)/N$。积分$C/T$得$S_X=\alpha T^3/3$、$S_G=\beta T$；能$E_X=E_0+\alpha T^4/4$、$E_G=E_0+\beta T^2/2$。形式自由能$F_X=E_0-\alpha T^4/12$、$F_G=E_0-\beta T^2/2$，相等$T_m^2=6\beta/\alpha$，假设熔化潜热$E_G-E_X=-\beta T_m^2=-6\beta^2/\alpha<0$，与所述结晶放热趋势不符。

官方解指出将冻结非遍历玻璃当作唯一平衡基态、套$S_G(0)=0$是可疑步骤。应理解为平衡第三定律的\textbf{适用前提}未满足，不能据此说平衡态第三定律被推翻。近似同$E_0$及将低温热容幂律外推到熔点也未经验证；玻璃转变本身含动力学，不能简单用两平衡相$\mu$相等预测。负潜热暴露了这些假设的共同不相容。
\end{solution}
\begin{exercise}[原题7(a)--(e)]\label{r:333-1-7}
求(a)$p=1/(2a),|x|<a$；(b)$p=\e^{-|x|/a}/(2a)$；(c)$p=a/[\pi(x^2+a^2)]$的特征函数、均值、方差；(d)$p=x\e^{-x^2/(2a^2)}/a^2,x\ge0$；(e)$p=\sqrt{2/\pi}x^2\e^{-x^2/(2a^2)}/a^3,x\ge0$仅求均值、方差。
\end{exercise}
\begin{solution}
用$\varphi(k)=\avg{\e^{\ii kx}}$，(a)$\sin(ak)/(ak)$，均0、方差$a^2/3$；(b)正负半轴指数积分给$1/(1+a^2k^2)$，均0、方差$2a^2$；(c)留数或已知Fourier积分$\e^{-a|k|}$，在0不可微，通常均值与方差都不存在，主值均0不等于期望。后两者用$\int_0^\infty x^m\e^{-x^2/(2a^2)}\dd x=\tfrac12(2a^2)^{(m+1)/2}\Gamma((m+1)/2)$：Rayleigh均$a\sqrt{\pi/2}$、方差$a^2(2-\pi/2)$；Maxwell均$2a\sqrt{2/\pi}$、方差$a^2(3-8/\pi)$。
\end{solution}
\begin{exercise}[原题8]\label{r:333-1-8}
任意概率密度$p(x)$均$\lambda$、有限方差$\sigma^2$，证明$\Pr(|X-\lambda|\ge n\sigma)\le1/n^2$。
\end{exercise}
\begin{solution}
$\sigma^2=\int(x-\lambda)^2p(x)\dd x\ge\int_{|x-\lambda|\ge n\sigma}(x-\lambda)^2p\dd x\ge n^2\sigma^2\Pr(|X-\lambda|\ge n\sigma)$，非零方差除去即得。零方差分布集中于均值，需以正偏离阈值处理，不能拿$0/0$论证。仅有限二阶矩保证，不要求Gaussian。
\end{solution}
\begin{exercise}[原题9(a),(b)]\label{r:333-1-9}
$n$个独立$r_\alpha\in[0,1]$同密度$p$，最大$X$。(a)密度；(b)均匀$p$的均值、方差和大$n$，讨论最优选择减少差异的条件。
\end{exercise}
\begin{solution}
$F_X(x)=F(x)^n$，$p_X=np(x)F(x)^{n-1}$，$F(x)=\int_0^xp$。均匀则$nx^{n-1}$，$\avg X=n/(n+1)$、$\avg{X^2}=n/(n+2)$，方差$n/[(n+1)^2(n+2)]\sim n^{-2}$，相对均值的散布收缩。$n(1-X)$收敛指数分布，仍有非Gaussian剩余波动。这解释此独立、有限上界模型的选择收缩，不是所有人群变异都由最优选择决定的结论。
\end{solution}
